% ==========================================
% SECCIÓN 14: REFERENCIAS BIBLIOGRÁFICAS
% ==========================================

\section{REFERENCIAS}

Las referencias bibliográficas se presentan bajo el estándar IEEE establecido por las normativas de titulación de la Facultad de Ingeniería y Ciencias Aplicadas (FICA) de la Universidad Central del Ecuador para la carrera de Ingeniería en Ciencias de la Computación, brindando el sustento teórico, normativo y tecnológico del proyecto integrador:

\begin{thebibliography}{99}

% Marco Legal y Normativo Institucional
\bibitem{loes}
Asamblea Nacional del Ecuador, \textit{Ley Orgánica de Educación Superior (LOES)}, Registro Oficial N.º 298, Quito, Ecuador, 2010.

\bibitem{lopdp_ecuador}
Asamblea Nacional del Ecuador, \textit{Ley Orgánica de Protección de Datos Personales}, Registro Oficial Suplemento N.º 459, Quito, Ecuador, 2021.

\bibitem{reglamento_uce}
Universidad Central del Ecuador, \textit{Reglamento de Régimen Académico Interno de la Universidad Central del Ecuador}, Res. RPC-SE-08-No.023-2022, Quito, Ecuador, 2022.

\bibitem{instructivo_titulacion}
Universidad Central del Ecuador, \textit{Instructivo General de Titulación de Grado de la Universidad Central del Ecuador}, Aprobado por el Honorable Consejo Universitario, Quito, Ecuador, 2023.

% Antecedentes Nacionales y Sistemas de Gestión Académica
\bibitem{romero_unesum_becas}
J. Romero y C. Medina, ``Diseño de un sistema informático para la automatización de la gestión documental en el área de becas de la Universidad Estatal del Sur de Manabí,'' Trabajo de titulación, Carrera de Tecnologías de la Información, Univ. Estatal del Sur de Manabí, Jipijapa, Ecuador, 2026.

\bibitem{barrionuevo_uta_tramites}
D. Barrionuevo, ``Sistema automatizado para la gestión de trámites académicos en la secretaría de la FISEI,'' Trabajo de titulación, Fac. Ing. en Sistemas, Electrónica e Industrial, Univ. Técnica de Ambato, Ambato, Ecuador, 2021.

\bibitem{senescyt_pusak}
Secretaría de Educación Superior, Ciencia, Tecnología e Innovación (SENESCYT), ``Plataforma PUSAK: Sistema de Gestión y Liquidación de Becas y Devengamiento Institucional,'' Quito, Ecuador, 2023. [En línea]. Disponible en: https://pusak.senescyt.gob.ec/

\bibitem{umet_expedientes}
Universidad Metropolitana del Ecuador, ``Modelo de Control Centralizado de Expedientes Docentes y Gestión de Información Académica,'' Guayaquil, Ecuador, 2022.

% Ingeniería de Software y Metodologías
\bibitem{pressman_eng}
R. S. Pressman y B. R. Maxim, \textit{Software Engineering: A Practitioner's Approach}, 9th ed. New York, NY, USA: McGraw-Hill Education, 2020.

\bibitem{sommerville_eng}
I. Sommerville, \textit{Software Engineering}, 10th ed. Boston, MA, USA: Pearson, 2016.

\bibitem{kendall_systems}
K. E. Kendall y J. E. Kendall, \textit{Systems Analysis and Design}, 10th ed. Upper Saddle River, NJ, USA: Pearson, 2019.

% Arquitectura de Software y Desarrollo Web Moderno
\bibitem{fowler}
M. Fowler, \textit{Patterns of Enterprise Application Architecture}. Boston, MA, USA: Addison-Wesley Professional, 2002.

\bibitem{martin_clean}
R. C. Martin, \textit{Clean Architecture: A Craftsman's Guide to Software Structure and Design}. Upper Saddle River, NJ, USA: Prentice Hall, 2017.

\bibitem{gamma_patterns}
E. Gamma, R. Helm, R. Johnson, y J. Vlissides, \textit{Design Patterns: Elements of Reusable Object-Oriented Software}. Reading, MA, USA: Addison-Wesley, 1994.

\bibitem{banks_react}
A. Banks y E. Porcello, \textit{Learning React: Modern Patterns for Developing React Applications}, 2nd ed. Sebastopol, CA, USA: O'Reilly Media, 2020.

\bibitem{vercel_nextjs}
Vercel Inc., \textit{Next.js 16 Documentation: Production-grade React applications that scale}, 2024. [En línea]. Disponible en: https://nextjs.org/docs

\bibitem{astro_docs}
Astro Technology Company, \textit{Astro Documentation: The web framework for content-driven websites}, 2024. [En línea]. Disponible en: https://docs.astro.build

\bibitem{tailwind_docs}
Tailwind Labs Inc., \textit{Tailwind CSS Documentation: Utility-first CSS framework}, 2024. [En línea]. Disponible en: https://tailwindcss.com/docs

\bibitem{express_docs}
StrongLoop y otros colaboradores, \textit{Express.js: Fast, unopinionated, minimalist web framework for Node.js}, 2024. [En línea]. Disponible en: https://expressjs.com/

\bibitem{prisma_docs}
Prisma Data, Inc., \textit{Prisma ORM: Next-generation ORM for Node.js and TypeScript}, 2024. [En línea]. Disponible en: https://www.prisma.io/docs

\bibitem{fumadocs_docs}
Fumadocs Core Team, \textit{Fumadocs: The Next.js framework for building documentation websites}, 2024. [En línea]. Disponible en: https://fumadocs.vercel.app/

% Sistemas de Bases de Datos Relacionales
\bibitem{elmasri}
R. Elmasri y S. Navathe, \textit{Fundamentos de Sistemas de Bases de Datos}, 7ma ed. Madrid, España: Pearson Educación, 2016.

\bibitem{silberschatz}
A. Silberschatz, H. F. Korth, y S. Sudarshan, \textit{Database System Concepts}, 7th ed. New York, NY, USA: McGraw-Hill, 2019.

\bibitem{postgresql_docs}
PostgreSQL Global Development Group, \textit{PostgreSQL 16.0 Documentation}, 2023. [En línea]. Disponible en: https://www.postgresql.org/docs/16/

% Seguridad Web, Autenticación y RBAC
\bibitem{ferraiolo_rbac}
D. F. Ferraiolo, D. R. Kuhn, y R. Chandramouli, \textit{Role-Based Access Control}, 2nd ed. Norwood, MA, USA: Artech House, 2007.

\bibitem{owasp_top10}
OWASP Foundation, \textit{OWASP Top 10 Web Application Security Risks}, Open Web Application Security Project, 2021. [En línea]. Disponible en: https://owasp.org/Top10/

\bibitem{rfc_jwt}
M. Jones, J. Bradley, y N. Sakimura, ``JSON Web Token (JWT),'' RFC 7519, Internet Engineering Task Force (IETF), May 2015.

% Integración de Sistemas y APIs RESTful
\bibitem{fielding_rest}
R. T. Fielding, ``Architectural Styles and the Design of Network-based Software Architectures,'' Tesis Doctoral, Dept. Inf. Comput. Sci., Univ. California, Irvine, CA, USA, 2000.

\bibitem{newman_microservices}
S. Newman, \textit{Building Microservices: Designing Fine-Grained Systems}, 2nd ed. Sebastopol, CA, USA: O'Reilly Media, 2021.

% Ingeniería de Requisitos y Evaluación de Software
\bibitem{ieee_requirements}
IEEE Computer Society, ``IEEE Standard for System and Software Engineering--Life Cycle Processes--Requirements Engineering,'' \textit{IEEE Std 29148-2018}, pp. 1--104, 2018.

\bibitem{nielsen_usability}
J. Nielsen, \textit{Usability Engineering}. San Diego, CA, USA: Academic Press, 1994.

\end{thebibliography}
